\documentclass[ngerman,11pt]{article}

\usepackage[utf8]{fontenc}
\usepackage[left=1.5cm, right=1.5cm, top=1cm, bottom=1cm]{geometry}
\usepackage{babel}
\usepackage{listings}
\usepackage{color}
\usepackage{hyperref}
\usepackage{csquotes}
\MakeOuterQuote{"}

\definecolor{dkgreen}{rgb}{0,0.6,0}
\definecolor{gray}{rgb}{0.5,0.5,0.5}
\definecolor{mauve}{rgb}{0.58,0,0.82}


\pagenumbering{gobble}
\renewcommand{\lstlistingname}{Schnipsel}

\lstset{frame=tb,
    language=Java,
    aboveskip=3mm,
    belowskip=3mm,
    showstringspaces=false,
    captionpos=b,
    columns=flexible,
    basicstyle={\small\ttfamily},
    numbers=left,
    numberstyle=\tiny\color{gray},
    keywordstyle=\color{blue},
    commentstyle=\color{dkgreen},
    stringstyle=\color{mauve},
    breaklines=true,
    breakatwhitespace=true,
    tabsize=3,
    inputencoding=utf8
}

% Packages
\usepackage{amsmath}

% Document
\begin{document}

    \setlength{\parindent}{0}

    {\huge \textbf{Informatik - Programmieren mit Java - Handout}}

    \vspace{1cm}


    \section{Arbeiten mit der Konsole}
    Die \textit{Konsole} ist ein Werkzeug zur Steuerung eines Computers.

    \vspace{0.5cm}

    \subsection{Konsole öffnen}
    Unter Windows können wir die Konsole unter anderem so öffnen:
    \begin{itemize}
        \item \lstinline{Windows + R} drücken, \lstinline{cmd} eingeben, \lstinline{Enter} drücken
        \item In der Windows-Suche (unten links) nach \lstinline{cmd} oder \lstinline{Eingabeaufforderung} suchen, auf \lstinline{Eingabeaufforderung} klcicken
    \end{itemize}

    \vspace{0.5cm}

    \subsection{Konsole bedienen}
    Öffnen wir die Konsole, sehen wir in etwa Folgendes:

    \begin{lstlisting}
C:\Users\Lernen>
    \end{lstlisting}

    Der Teil vor dem \lstinline{'>'} ist der Pfad, "an dem wir uns befinden" - das sogenannte \textit{Arbeitsverzeichnis}. \\

    Wir haben 2 Konsolen-Befehle kennengelernt:
    \begin{itemize}
        \item \lstinline{cd} \textit{(\textbf{c}hange \textbf{d}irectory)} - wechselt das Arbeitsverzeichnis. \\
        Anwendung:
        \begin{lstlisting}
cd Unterordner
cd C:\Users\Lernen\Desktop\MeinOrdner
cd ..
        \end{lstlisting}
        (Zeile 1: Wechsle zu einem \textit{Unterodner} im aktuellen Ordner) \\
        (Zeile 2: Wechsle zum angegebenen Pfad) \\
        (Zeile 3: Wechsle zum \textit{Überordner}) \\

        \item \lstinline{dir} \textit{(directory)} - zeigt alle Dateien und Ordner im aktuellen Arbeitsverzeichnis an
    \end{itemize}

    \vspace{0.5cm}


    \section{Java installieren}
    Bevor wir mit Java arbeiten können, müssen wir das \textit{JDK} (Java-Development-Kit) installieren.
    Am einfachsten geht das über die Konsole mit diesem Befehl:

    \begin{lstlisting}
winget install EclipseAdoptium.Temurin.17.JDK
    \end{lstlisting}

    Schließen wir nach erfolgter Installation die Konsole und öffnen sie erneut, können wir mit dem
    Befehl \lstinline{java} prüfen, ob die Installation erfolgreich war.
    War Sie es nicht, sehen wir einen Fehler, in etwa "Den Befehl gibt es nicht" -
    ansonsten einen längeren Hilfetext von Java selbst.


    \section{Java - absolute Grundlagen}

    \vspace{0.5cm}

    \subsection{Java-Programm ausführen}
    Wir haben die Grundlagen betrachtet, wie man mit der Programmiersprache Java arbeitet.
    Noch einmal zusammengefasst:

    \begin{itemize}
        \item Wir schreiben \textit{Programmcode} (oder einfach \textit{Code}) in einer Datei, z.B. \lstinline{TollesProgramm.java}
        \item Wir \textit{kompilieren} das Programm mit \lstinline{javac TollesProgramm.java}
        \item Wir führen es aus mit \lstinline{java TollesProgramm}
    \end{itemize}

    Der Befehl \lstinline{javac} erzeugt eine Datei \lstinline{TollesProgramm.class}, wenn alles funktioniert hat.
    Ansonsten bekommen wir eine \textit{Fehlermeldung}.
    Letztere beschreibt mal mehr mal weniger genau, wo wir einen Fehler gemacht haben. \\

    \textit{Beachte}: Damit die Befehle \lstinline{javac} und \lstinline{java} funktionieren, müssen wir im korrekten Arbeitsverzeichnis sein
    (nämlich dort, wo sich unsere Datei \lstinline{TollesProgramm.java} befindet).

    \vspace{0.5cm}

    \subsection{Minimales Java Programm}
    Zuletzt haben wir ein minimales Java-Programm betrachtet.
    Dieses sieht wie folgt aus:

    \begin{lstlisting}[label={lst:minimal_java},caption=Minimales Java Programm]
public class Test {
    public static void main(String[] args) {
    }
}
    \end{lstlisting}

    Dieses Programm tut gar nichts.
    Diese 4 Zeilen stellen das "Grundgerüst" eines Java-Programms dar. \\

    Erste neue Kleinigkeit: mit \lstinline{//} (doppel-slash) können wir \textit{Kommentare} in unserem
    Code schreiben.
    Alles innerhalb einer Zeile, was nach \lstinline{//} kommt, wird zu einem Kommentar.
    Ein Kommentar wird von Java komplett ignoriert, und taucht im kompilierten Code (der \lstinline{.class}-Datei) nicht auf. \\

    Wir nutzen Kommentare meistens, um anderen Lesern des Codes (oder unserem zukünftigen Ich) zu erläutern, WARUM
    bestimmte Dinge im Code so sind wie sie sind (WAS passiert brauchen wir nicht zu erklären, das sagt ja der Programmcode). \\

    Die folgenden Zeilen sind kompiliert also exakt das selbe Programm wie in Schnipsel 1:

    \begin{lstlisting}[caption=Kommentare]
// Ich bin nur ein Kommentar und habe sonst keine Funktion
public class Test { // Ich bin auch ein Kommentar (der Teil vor mir aber nicht)
    public static void main(String[] args) {
        // Ich bin noch ein Kommentar
        // Wieder ein Kommentar
    }
    // Ich stehe hier nur so
}
    \end{lstlisting}

    \vspace{0.5cm}

    \subsection{Der main-Block}
    Wir können noch nicht jedes Detail des Grundgerüsts klären.
    Es reicht aber zunächst zu wissen, dass unser Programmcode in den inneren Block, den "main-Block" gehört:

    \begin{lstlisting}[caption=main-Block]
public class Test {
    public static void main(String[] args) {
        // Das hier ist der main-Block unseres Programms.
        // Hier schreiben wir (vorerst) unseren gesamten Programmcode rein
    }
}
    \end{lstlisting}


    \section{Java - Grundlagen}

    Die Befehle, die wir in den main-Block schreiben, werden nacheinander von oben nach unten abgearbeitet. \\

    Wir halten einmal die ersten wichtigen Java Befehle fest.
    Jeder Befehl wird mit einem \lstinline{;} \textit{(Semikolon)} beendet.

    \vspace{0.5cm}
    \subsection{Ausgabe}

    Mit \lstinline{System.out.println} können wir unser Programm etwas auf der Konsole ausgeben lassen.
    Auszugebenden Text setzen wir dabei in Anführungszeichen (Erklärung im Abschnitt "Strings").
    Beispiel:

    \begin{lstlisting}[label={lst:sysout},caption=Ausgabe auf der Konsole]
public class Test {
    public static void main(String[] args) {
        System.out.println("Ich bin ein interessanter Text");
        System.out.println("42");

        // Die nächste Zeile produziert die selbe Ausgabe wie die vorherige, die Zahl wird einfach zu Text umgewandelt:
        System.out.println(42);
    }
}
    \end{lstlisting}

    \newpage
    \vspace{0.5cm}
    \subsection{Variablen}

    Eine Variable ist ein \textit{Platzhalter}.
    Eine Variable hat stets einen
    \begin{itemize}
        \item \textit{Datentyp} (oder einfach Typ),
        \item einen \textit{Namen},
        \item und einen \textit{Wert}.
    \end{itemize}

    Der Datentyp gibt an WAS FÜR EINE ART von Werten in der Variable gespeichert werden kann.
    Den Namen nutzen wir, um auf den Wert der Variablen zuzugreifen oder allgemein "etwas mit ihr zu tun".\\

    Beispiel:

    \begin{lstlisting}[caption=Variable]
public class Test {
    public static void main(String[] args) {
        int anzahlBananen = 17;
        System.out.println(anzahlBananen);
    }
}
    \end{lstlisting}

    Hier erstellen wir eine Variable vom Typ \lstinline{int} (das steht für \textit{Integer}, englisch für Ganzzahl)
    und weisen ihr den Wert $17$ zu. \\

    Je nach Datentyp können wir verschiedene Dinge mit Variablen anstellen. \\

    Zahlen können wir zum Beispiel mit \lstinline{+} addieren:

    \begin{lstlisting}[caption=Summe zweier Ganzzahlen]
public class Test {
    public static void main(String[] args) {
        int a = 15;
        int b = 21;
        int summe = a + b;
        System.out.println(summe);
    }
}
    \end{lstlisting}

    und genauso subtrahieren (\lstinline{-}), multiplizieren (\lstinline{*}) oder dividieren (\lstinline{/}).

    Java kann auch komplexere Terme interpretieren:

\begin{lstlisting}[caption=binomische Formel]
public class Test {
    public static void main(String[] args) {
        int a = 4;
        int b = 5;
        int ergebnis1 = (a + b) * (a + b);
        int ergebnis2 = a * a + 2 * b * a + b * b;

        // In Andenken an die binomische Formel hoffen wir mal dass hier zweimal das selbe rauskommt:
        System.out.println(ergebnis1);
        System.out.println(ergebnis2);
    }
}
\end{lstlisting}


    \vspace{0.5cm}
    \subsection{Einer Variablen einen neuen Wert zuweisen}

    Mit dem Befehl \lstinline{int x = 5} tun wir eigentlich 2 Dinge auf einmal:

    \begin{itemize}
        \item Wir \textit{deklarieren} die Variable \lstinline{x} vom Typ \lstinline{int}
        (wir "sagen dass es sie gibt", englisch "to declare"),
        \item und wir weisen \lstinline{x} den Wert $5$ zu.
    \end{itemize}

    Wenn wir den Wert einer Variablen nachträglich ändern wollen, nutzen wir dieselbe
    schreibweise, aber ohne den Typ zu wiederholen:

    \begin{lstlisting}[caption=Wert einer Variablen ändern]
public class Test {
    public static void main(String[] args) {
        int x = 5;
        System.out.println(x);
        x = -3;
        System.out.println(x);
    }
}
    \end{lstlisting}

    Variablennamen müssen eindeutig sein (innerhalb eines "scopes", dazu später mehr).
    Schreiben wir \lstinline{int x = 5} und später \lstinline{int x = -3}, denkt der Compiler,
    dass wir zwei verschiedene Variablen mit dem Namen \lstinline{x} haben wollen - das führt zu einem Fehler.

    \subsection{Weitere Datentypen}

    Weitere Datentypen sind:
    \begin{itemize}
        \item \lstinline{double} (für Kommazahlen, z.B. \lstinline{double x = 0.5543})
        \item \lstinline{boolean} (für Wahrheitswerte, z.B. \lstinline{boolean T = true})
        \item \lstinline{String} (für Zeichenketten, z.B. \lstinline{String text = "Bananenkuchen"})
    \end{itemize}

    \textit{Beachte}: \lstinline{String} gehört eigentlich zu einer anderen Kategorie als \lstinline{int}, \lstinline{double}, \lstinline{boolean}.
    Letztere sind sogenannte \textit{primitive Datentypen}, \lstinline{String} dagegen \textbf{nicht} - aber dazu später mehr. \\

    Anmerkung: Wir können auch \lstinline{String}'s addieren:

    \begin{lstlisting}[caption=Strings addieren]
public class Test {
    public static void main(String[] args) {
        String teil1 = "Ciao";
        String teil2 = "Kakao";
        System.out.println(teil1 + " " + teil2);
    }
}
    \end{lstlisting}

    Wir können Strings sogar mit beliebigen anderen Variablen addieren:

    \begin{lstlisting}[caption=String + Zahl]
public class Test {
    public static void main(String[] args) {
        int x = 18;
        System.out.println("Hier liegen " + x  + " Maracujas");
    }
}
    \end{lstlisting}

    Die Summanden, die nicht vom Typ \lstinline{String} sind, werden dabei zu Strings umgewandelt.

    \subsection{Bedingungen und Kontrollstrukturen}

    Unser Programm wird zwar von oben nach unten, Befehl für Befehl abgearbeitet - dennoch wollen
    wir manchmal dass es sich nicht immer exakt gleich verhält.

    Eine Art dies zu tun, ist, den Programmfluss an \textit{Bedingungen} zu knüpfen - einfach ausgedrückt:
    "Wenn das, dann\dots\ sonst etwas anderes". \\

    Um zu überprüfen "ob etwas gilt" brauchen wir Wahrheitswerte (englisch: boolean).
    Eine Variable vom typ \lstinline{boolean} kann nur genau zwei Werte haben: \lstinline{true} oder \lstinline{false}.
    \textit{True} ist englisch für "wahr", \textit{False} für "falsch".\\

    Wir können so einen Wahrheitswert zum Beispiel durch den Vergleich von Zahlenwerten vergleichen:

    \begin{lstlisting}[caption=Wahrheitswerte]
public class Test {
    public static void main(String[] args) {
        int x = 4;
        int y = 23494;

        boolean xKleinerY = x < y;
        boolean yKleinerX = y < x;
        boolean xGroesserY = x > y;
        boolean yGroesserX = y > x;
        boolean xGleichY = x == y;
        boolean xUngleichY = x != y;

        System.out.println(xKleinerY);
        System.out.println(yKleinerX);
        // ...
    }
}
    \end{lstlisting}


    Wahrheitswerte können wir nutzen, um bestimmte Teile des Programms optional ausführen zu lassen - nämlich
    nur dann wenn eine Bedingung wahr ist. \\

    Das geht wie folgt:

    \begin{lstlisting}[caption=Variable]
public class Test {
    public static void main(String[] args) {
        int x = 5;
        double y = 3.4;

        // Der Code innerhalb des if-Blocks wird nur ausgefuehrt, wenn das Resultat von "x > y" true ist:
        if(x > y) {
            System.out.println("x ist groesser als y!!!!");
        }
    }
}
    \end{lstlisting}

    Diese Struktur nennt man \textit{if-Abfragge} oder \textit{if-Statement} - \textit{if} bedeutet auf englisch "wenn". \\

    Ein \lstinline{if} können wir zudem mit einem \lstinline{else} (englisch für sonst) kombinieren:

    \begin{lstlisting}[caption=Variable]
public class Test {
    public static void main(String[] args) {
        int x = 5;
        double y = 3.4;

        // Der Code innerhalb des if-Blocks wird nur ausgefuehrt, wenn das Resultat von "x > y" true ist:
        if(x > y) {
            System.out.println("x ist groesser als y!!!!");
        } else {
            // Der code in diesem else Block wird nur "sonst" ausgeführt, also wenn die Bedingung im if(...) NICHT wahr ist
            System.out.println("x ist anscheinend nicht groesser als y");
        }
    }
}
    \end{lstlisting}

    Konkretes Beispiel: Wir wollen entscheiden, ob ein Kreis mit dem angegebenen Radius größer oder kleiner als ein Grenzwert ist.

    \begin{lstlisting}[caption=Kreisfläche]
public class Test {
    public static void main(String[] args) {
        double PI = 3.141

        double radius = 13;
        double grenzwert = 30;

        double flaeche = PI * radius * radius;

        if(flaeche > grenzwert) {
            System.out.println("Die Flaeche ist " + flaeche + ", groesser als der Grenzwert!");
        } else {
            System.out.println("Die Flaeche ist " + flaeche + ", das ist ok.");
        }
    }
}
    \end{lstlisting}

\end{document}
